\chapter{Espacios de Hardy}

\transclude{norma-lp}

\transclude{esp-lp}

\begin{prop}\label{prop:fn-holomorfa-esencialmente-acotada-imp-acotada}
    Sea $f \in \mathcal{H}\left(\mathbb{D}\right)$ tal que $f$ está esencialmente acotada en $\mathbb{D}$, es decir, $f \in \mathcal{L}^{\infty} \left(\mathbb{D} \right)$. Entonces, $f$ es acotada.
\end{prop}
\begin{dem}
    Sea $M \defeq \norm{f}_{\mathcal{L}^{\infty}(\mathbb{D})} < \infty$. Fijamos $z_0 \in \mathbb{D}$ y tomamos $r_0 > 0$ tal que $\overline{D(z_0,r_0)} \subset \mathbb{D}$. Para todo $\rho \in (0,r_0)$, aplicando la \exhyperref[teo:valor-medio-fn-holomorfa]{teo-valor-medio-fn-holomorfa}{propiedad de valor medio para funciones holomorfas} en $D(z_0,\rho)$,
    \[
    f(z_0) = \frac{1}{2\pi} \int_0^{2\pi} f\left(z_0+\rho e^{it}\right) \, \odif{t}.
    \]
    Multiplicando por $\sfrac{2\rho}{r_0^2}$ e integrando en $\rho \in (0,r_0)$, obtenemos
    \[
    f(z_0) = \frac{2}{r_0^2} \int_0^{r_0} \left(\frac{1}{2\pi} \int_0^{2\pi} f\left(z_0+\rho e^{it}\right) \, \odif{t}\right) \rho \, \odif{\rho}.
    \]
    \[
    \begin{aligned}
        \therefore \quad\abs{f(z_0)} &\leq \frac{2}{r_0^2} \int_0^{r_0} \left(\frac{1}{2\pi} \int_0^{2\pi} \abs{f\left(z_0+\rho e^{it}\right)} \, \odif{t}\right) \rho \, \odif{\rho} \\
        &= \frac{1}{\pi r_0^2} \int_{D(z_0,r_0)} \abs{f(z)} \, \odif{A}(z) \leq \frac{1}{\pi r_0^2} \int_{D(z_0,r_0)} M \, \odif{A}(z) = M.
    \end{aligned}
    \]
    Como $z_0$ era arbitrario, $\forall z \in \mathbb{D} : \abs{f(z)} \leq M$. Luego $f$ es acotada en $\mathbb{D}$.
\end{dem}